% Draft 2026-11-03, official setup. No claim here about memory gains until results/official/analysis.json exists; \todo{} marks them.
\section{Discussion}
\label{sec:discussion}

\textbf{Links behind a score.}
Our audit finds a risk at each model-independent link in Table~\ref{tab:chain}.
For 8 of 15 tasks no corpus run resembles the task's prompt; the reference solution scores 0 on 4 tasks, as the maintainers document, yet without memory our agent passed two of them in the counted control (cookie store in 1 of 3 attempts, EnterpriseOps CSM in 2 of 3); the same entry file is 21--100\% visible depending on the harness; and a slow model's calls are cut by the network proxy or abandoned by the agent and scored as failed attempts.
A score difference between two memory systems mixes memory quality with corpus coverage, harness loading and grader limits, and the leaderboard publishes too little to separate them.
With this model on this pipeline, generated memory gave no measurable gain over no memory, and the interval excludes mean gains above 0.14 per task.

\textbf{What is already documented.} The maintainers list the four zero reference scores~\cite{frontiermemory-code}, report a mean and standard deviation across 3 runs~\cite{frontiermemory}, updated the corpus to v1.1 to reduce reward hacking~\cite{frontiermemory-changelog}, and plan an online setting~\cite{frontiermemory}. The rules fix the task model to Opus 5.5 or GPT-6 Sol~\cite{frontiermemory-code}, so our local-model results lie outside what the leaderboard supports.

\textbf{Recommendations for the benchmark.}
\begin{itemize}
  \item Publish the no-memory rows the rules already mention, with per-task results and all attempts, so readers can test score gaps.
  \item Report per-task corpus coverage, and the judge model for judged tasks.
  \item Score a trial whose LLM judge failed as missing, and report it as an infrastructure failure.
  \item Pin build dependencies, or ship prebuilt images, so a task cannot change under a published score.
  \item Make the submission check do what the rules say: compare the corpus digest and the task and attempt counts.
  \item Report how much of each memory the harness shows at start, and how often agents open the rest.
  \item Treat a model call cut by the network proxy or abandoned by the agent's own timeout as an infrastructure failure, and document the proxy's 15-second first-byte limit (the proxy belongs to Harbor, upstream of the benchmark).
\end{itemize}

\textbf{Threats to validity.}
\begin{itemize}
  \item \emph{Model.} One local 27B model generates memory and runs tasks. Our findings about gains from memory, generation quality and agent behavior apply to it only; findings about corpus, graders, reporting and caps do not depend on it.
  \item \emph{Harness.} We run Hermes, a published row, because the default generator and loader are not public. Results may differ for Letta Dream and Letta Code. On a second machine, Codex with the same run-1 memory passed 7 of 45 attempts against 6 without (a separate exploratory study).
  \item \emph{Serving.} One model slot serves each trial, so side calls wait behind long ones; the relay and longer call timeouts, the same in every arm, remove the failures this caused. Hermes sends \emph{xhigh} effort for an unrecognized model, and the server barely acts on any setting: on a fixed probe, reasoning length was similar for low, \emph{xhigh}, max and no effort at all, so ``highest effort'' has little effect on this model; a second GPU type agreed. The proxy and timeout findings need a model slow enough to hit them.
  \item \emph{Deviations.} All are listed in \S\ref{sec:method}.
  \item \emph{Order.} The control ran before the memory runs rather than interleaved; the environment and model are frozen, and the remote judge can change between runs.
  \item \emph{Corpus coverage.} Prompt similarity misses runs on the same repository with different prompts, so a task with no matching run is not proven unlearnable.
\end{itemize}
